\section{Recovering the prescribed downward theta test}

% BEGIN Q87F1
\paragraph{Q87F1. The good projector commutes with the prescribed lift.}
Fix an original realization with the level, coefficient, Weil and
theta data in $\mathrm{Given}$ and $\mathrm{FullP}$, a row, and a twist.
Write $V$ for the definite source of the downward lift, $W$ for its
target, $C$ for the source coefficient space, $v_P$ for the chosen
Schwartz vector and $I_y(f,\ell,g)$ for the integral in Y86N2.
On the good source subspace one has
\[
             e_W^{\rm inc}\,\mathrm{thG}(y)=\mathrm{thF}(y).
\]
Here $e_W^{\rm inc}$ includes the native lower space in the space of
automorphic functions with coefficients. To check the equality, let
$e_0=\sum_b c_b b$ be the good idempotent in the common finite group
algebra. Lift each $b$ on both sides to finite parahoric elements.
Surjectivity on $V$ follows from the level realization; on $W$ lift
the invertible residue blocks in their frames and complete by the
identity. The partners above the split primes are determined by the
Hermitian adjoint, and the designated primes avoid their conjugates.

The normalized action $a_V(b)$ on $y$ is right translation by the
lift divided by its flag character. It is independent of the lift
because the ratio of two lifts lies in the principal-type kernel and
acts by its flag character. The same holds on $W$. For each lowest
target tensor $x$, changing the integration variable by the $V$-lift
gives
\[
 \mathrm{thF}(a_V(b)y)(x)
       =a_W(b)\mathrm{thF}(y)(x).
\]
Indeed the paired changes of frame preserve the matrix residue test
of $v_P$ up to the ratio of the two flag characters. At a type
prime both changes carry the distinguished residue frame into the
same frame; at the other level primes they have principal residue.
The Weil factor on determinants is the level-character ratio by
the defining law of the chosen step vector. The character
corrections in the two translations cancel it.
The finite translations commute with the real operators making
the lowest coefficient polynomial from $v_P$. The definite quotient
has right-invariant measure.

The good hypothesis fixes every scalar component of $y$ under
$e_0$. Linear functionals separate the finite-dimensional space $C$,
so the extended source functions of $\sum_b c_b a_V(b)y$ and $y$
agree. Summing the displayed equality proves that the full lift is
already fixed by $e_0$ on $W$ on each lowest tensor. The definition
of $\mathrm{thG}$ applies precisely this projector to the full lift.
The lowest tensors generate the target coefficient module; equality
of their evaluations extends to the two automorphic maps as in
E77A2 and E77A5. This proves $\mathrm{GoodMem}$.
% END Q87F1

% BEGIN Q87F2
\paragraph{Q87F2. Recovering every polynomial coefficient at infinity.}
Let $\Phi(P)$ be the member's Fock polynomial $P$ applied to $v_P$,
with the finite vector held fixed. If $\mathrm{thF}(y)=0$, then
$I_y(\Phi(P),\ell,g)=0$ for every $P,\ell,g$.

For a pure tensor $x$ of lowest $W$-types, the prescribed polynomial
is $\mathrm{io}(x,w)$, linear in $w\in C$. At fixed $g$ the bilinear
form $(w,\ell)\mapsto I_y(\Phi(\mathrm{io}(x,w)),\ell,g)$ gives an
endomorphism $A_{x,g}$ of $C$ by evaluation against $C^*$.
Right translation in the compact real group of $V$ gives
$A_{x,g}\rho(k)=\rho(k)A_{x,g}$: the source coefficient changes by
$\rho(k)^{-1}$ and the polynomial transforms by $\rho(k)$.
The coefficient representation is irreducible, so $A_{x,g}$ is
scalar. In a basis $w_j$ of $C$ with dual $w_j^*$, its trace is
\[
 \sum_j I_y(\Phi(\mathrm{io}(x,w_j)),w_j^*,g)
                     = \mathrm{thF}(y)(x)(g).
\]
This is the expansion of the vector-valued integral in that basis.
The scalar is zero since $\dim C>0$. Thus the vanishing holds for
every individual lowest coefficient and, by linearity, for every
tensor $x$.

Take the common kernel in polynomial space of the functionals with
$\ell$ and $g$ variable. It is stable under the compact source group
by the same change of variables. It is stable under the real target
Lie algebra because translating a test by the target changes $g$ by
right multiplication. Differentiate this equality along a real
one-parameter subgroup. The differentiated polynomial is its usual
quadratic Fock Lie operator including the splitting-character scalar.
The derivative passes under the integral: on compact representatives
of the definite quotient the translated Schwartz seminorms and the
first derivatives are uniformly bounded, while the source
coefficient is integrable. TAE3--6 give this argument in the
quasiperiodic realization. Complex linearity gives the whole
complexified target Lie algebra.

In each bounded polynomial degree the compact source representation
is completely reducible. The pairing with $C^*$ vanishes on every
source type different from $C$, by equivariance and Schur's lemma.
On its $C$-isotypic part, the compact dual-pair theorem says that the
multiplicity space $\operatorname{Hom}_{U(V)_\infty}(C,\mathcal P)$
is irreducible under the complexified target Lie algebra and its
maximal compact group (Howe, \emph{Transcending classical invariant
theory}, \S3, (3.9)). Here the pairs are the actual
$U(\dim V)\times U(p_\tau,q_\tau)$, one per real place, with the
splittings of JIO13--18. The finite product has the external tensor
product of the local multiplicity modules. The source coefficient is
irreducible for this product, and the maximal compact target factors
are connected. On the polynomial module their locally finite action
is generated on each compact orbit span by the differentiated
action. The maps $w\mapsto\mathrm{io}(x,w)$ are intertwiners,
and at least one is nonzero by JIO11--12 and E77A4. Those with image
in the common kernel form a nonzero stable submodule of the
multiplicity space. They therefore exhaust it. This proves the
assertion for every polynomial.
% END Q87F2

% BEGIN Q87F3
\paragraph{Q87F3. Fixed finite annihilator and Sobolev approximation.}
Let $f$ be a member Schwartz test such that every finite Heisenberg
sum killing $v_P$ kills $f$. Then $f$ is a limit of the vectors
$\Phi(P)$ in every one of the source's real Heisenberg Sobolev
seminorms, on one common finite level.

Use a rational polarization of the underlying symplectic space.
In its Schr\"odinger coordinates a test at fixed finite level is a
finite sum of Schwartz functions of the real coordinates times
finite residue functions. One can see this directly in the original
theta model by taking Fourier coefficients along one rational
Lagrangian modulo its rational points. Quasiperiodicity along the
other rational Lagrangian gives all coefficients as translates of the
coefficient at zero, and Fourier inversion recovers the given
theta function. Integrating by parts in the compact real Fourier
variables gives rapid decay of those coefficients with all derivatives.
At finite places invariance under an open lattice of translations and
under an open lattice of multiplications by additive characters gives,
respectively, periodicity and bounded support in the annihilator of
the latter lattice. Passing to a smaller common finite level works
for both $f$ and $v_P$. Its finite factors lie in the functions on a
single finite quotient $B/A$ of compact open lattices. Outside a
finite set the integral lattices are self-dual and the corresponding
factor has dimension one.

The vector $v_P$ has a single finite factor $v_f$, up to an overall
nonzero scalar. Indeed its real vacuum equations give the single
Gaussian line at infinity, and its finite formula consists of the
distinguished lattice vacuum at the ordinary places and the product
of the residue functions at the type and twist places. On functions
on $B/A$, additions by $B/A$ and multiplications by its dual
characters span the full matrix algebra: Fourier-average the
multiplications to extract an individual coordinate and translate it
to each other coordinate. Lift the translations and characters to
finite Heisenberg operators on the adeles. Their central phases are
scalars and all these operators preserve this finite model.
In particular there is a finite Heisenberg sum $p$ whose action
on it projects onto $\mathbb C v_f$ and fixes $v_f$.
Now $(p-1)v_P=0$, whence $(p-1)f=0$. Thus $f=f_\infty\otimes v_f$
in these coordinates. This argument permits an arbitrary finite
tensor for $f$ before the projection.

Expand $f_\infty$ in the Hermite basis of the positive complex
structure determined by the member's real frames. The oscillator
has eigenvalues $2|\alpha|+d_0$ after a linear rescaling, with $d_0$
the real configuration dimension. Every oscillator power applied
to $f_\infty$ belongs to $L^2$. Parseval therefore bounds the
squared coefficients with every polynomial weight. Each Heisenberg
derivative is a sum of raising and lowering operators with
coefficients of size $O((1+|\alpha|)^{1/2})$.
Truncation by total degree consequently converges in the $L^2$ norm
of every finite word in these derivatives. The same holds for
their finite sum defining $\mathrm{sobO}$, with the covolume-one
measure, by Plancherel for the preceding Fourier model. These
truncations are exactly polynomials applied to the prescribed
Gaussian times $v_f$, hence have the form $\Phi(P)$.

For fixed $y,\ell,g$ the theta integral is continuous on this
level. On compact representatives of the definite quotient the
value of the theta kernel is bounded by $C_g\,\mathrm{sobO}_d$
for some finite $d$. This follows either from TAE4--6 or by lattice
summation with a Schwartz weight of exponent greater than the real
dimension on the common bounded finite support. Compactness makes
the bound uniform in the source variable, and
$\int |\ell(\operatorname{ext}_V(y))|$ is finite. Applying the
bound to the difference of two tests proves continuity. Q87F2 now
proves $\mathrm{GenInf}$ for the member.
% END Q87F3

% BEGIN Q87F4
\paragraph{Q87F4. Removing the principal-level differences and the bad averages.}
Write $L=\mathrm{datVt}(i,t)$ and
$J=\mathrm{levG}(V,\beta_V,L)$, with character
$\lambda=\mathrm{levChi}(V,\beta_V,L)$. Here the quotient level
of $L$ is trivial, so $\mathrm{flag1L}(L)=L$. In particular at a
type prime the indicated residue frame is fixed pointwise, whereas
the twist primes retain their parahoric character. Every theta integral of $y$ on an arbitrary test $f$
equals its integral on a test of type $\lambda^{-1}$.
Indeed, if $k\in J$, change the variable $h$ to $hk$ in the
defining integral. The source component has right character
$\lambda$, giving
\[
 I_y(\omega(k,1)f,\ell,g)=\lambda(k)^{-1}I_y(f,\ell,g).
\]
Averaging $\lambda(k)\omega(k,1)f$ over $J$ proves the assertion.
The test and character are smooth at finite places, so the
probability average is a finite sum.

On this type space a lift $k_b$ of a finite quotient element $b$
acts on tests by
$\lambda_{\rm flag}(k_b)^{-1}\omega(k_b^{-1},1)$.
Changing the lift by a principal element leaves this operator
unchanged. Transferring it across the integral yields exactly the
corrected action $\lambda_{\rm flag}(k_b)^{-1}R(k_b)$ on
$\ell(\operatorname{ext}_V y)$. Suppose $y$ has the two vanishing
mirabolic averages, for both choices of character exponent, as in
N86V1. Then every average of these test operators with the source's
weights integrates to zero. In the span of the finite orbit of a
given type test, take an invariant positive Hermitian inner product
by averaging any such product. Every weighted subgroup average,
divided by the subgroup order, is an orthogonal character projector
(the lift action on tests has the opposite order, which preserves
the subgroup average). The orthogonal complement of the span of
all their images is their common kernel. Decomposing the type test
in this finite-dimensional space leaves a test killed by all the
averages and having the same integrals. We call such a test clean.
Taking all defining vectors, both subgroups and both exponents makes
the collection invariant under conjugation. All the constructions
use only finite source operators and commute with target translations.
% END Q87F4

% BEGIN Q87F5
\paragraph{Q87F5. Generation at ordinary finite places.}
In the local factor of a clean type test at an ordinary place,
the source level fixes the test and the prescribed vector is the
standard lattice vacuum. The source-invariant test belongs to
the algebraic span of target translates of that vacuum.

At a split place use a half of the split symplectic space as
Schr\"odinger coordinates. Put $K=F_v$, $O=\mathcal O_v$, $n=r+t$,
and identify the half, using the member's self-dual frames, with
$M_{n,r}(K)$. Source integral matrices act on the left, target
matrices on the right, up to the usual common nonzero Weil scalars.
The scalar on the source integral group is one, as that group
fixes the integral vacuum $\mathbf1_{M_{n,r}(O)}$. A locally
constant compactly supported invariant $H$ has support in
$d^{-1}M_{n,r}(O)$ and period $dM_{n,r}(O)$ for some nonzero
$d\in O$. For $X$ in this support bound, form the row span
$L_X$ of $dX$ modulo $d^2$ in $(O/d^2)^r$.
Two matrices with the same row span are related by an invertible
row operation modulo $d^2$. In fact over the local principal
quotient any generating tuple of a given module can be changed
into a minimal generating tuple followed by zeros: select a basis
in its residue quotient, apply Nakayama and subtract linear
combinations from the extra rows. Two minimal tuples are related by
a matrix invertible on reduction, which lifts to an integral
invertible matrix. Periodicity now shows $H$ depends only on $L_X$.
For any submodule $L$ of $(O/d^2)^r$, the condition $L_X\subseteq L$,
together with the support bound, says that all the rows of $X$
belong to the full lattice $d^{-1}\widetilde L$, where
$\widetilde L$ is the inverse image in $O^r$. Its indicator
is a right $\mathrm{GL}_r(K)$ translate of the standard vacuum.
Subtracting the indicators of the proper submodule strata
recursively gives every exact-row-span indicator. The submodule
poset is finite. This proves generation, also for functions taking
values in a fixed auxiliary vector space.

At a nonsplit place the host is unramified and places over $2$ split,
so the local residue characteristic is odd. The source and target
lattices are self-dual. In self-dual coordinates the trace pairing
has the primitive conductor prescribed by the source's standard
additive character (absorb the inverse different in the frames).
The two Hermitian Gram matrices are integral and unimodular;
a trace-zero unit gives the skew normalization. The vacuum is
fixed by both integral unitary groups, with their lattice splittings.
Apply the odd-residue-characteristic lattice-model theorem of
M\oe glin--Vign\'eras--Waldspurger, \emph{Correspondances de Howe
sur un corps $p$-adique}, theorem on p.~103, with the full self-dual
source lattice $L_1$ and the self-dual target lattice $L_2$.
In the notation of that theorem $L=L_1=L_1^\perp$, hence
$B(L)=L^\perp\otimes L_2=L_1\otimes L_2$. A function in the
lattice model supported on this one lattice coset is determined by
its value at zero. This gives the vacuum line. The condition
$(u-1)L^\perp\subseteq L$ holds for every $u$ stabilizing $L_1$,
and the odd-characteristic lattice splitting fixes the vacuum.
Thus the group $J_1(L)$ here is the full integral source group. The theorem identifies
its invariants with the target Hecke algebra applied to this line.
An operator with compact support on a smooth vector is a finite
linear combination of group translates. This is precisely the
source-fixed branch of $\mathrm{HoweUnram}$. One can equally apply
the pinned proof of this exact lattice theorem. Tensoring with
auxiliary coefficients leaves the assertion intact, by taking a
basis of the span of the finitely many coefficients.
% END Q87F5

% BEGIN Q87F6
\paragraph{Q87F6. The frame and determinant tests at the two kinds of level prime.}
At a designated type prime use the conjugate split half and the
member's ordered source and target frames. In the resulting
$M_{r+t,r}(K)$ coordinates put $e=\binom{1_r}{0}$ and
$B=e+\varpi M_{r+t,r}(O)$, where $\varpi$ is a uniformizer.
The prescribed test on this half is $\mathbf1_B$ up to a nonzero
scalar. In these coordinates the finite restriction on a clean
type test $H$ has two consequences:
\begin{enumerate}
\item $H$ is invariant under integral left multiplication fixing
the columns of $e$ pointwise modulo $\varpi$.
\item For $0\ne a\in(k^r)^*$ and $D_a\subset\mathrm{GL}_r(k)$
the subgroup fixing $\ker a$ pointwise, one has
$\sum_{b\in D_a}H(\operatorname{diag}(\widetilde b,1)X)=0$.
\end{enumerate}
Here $k=O/\varpi$ and $\widetilde b$ is any integral lift. For the
first assertion the lift lies in the principal parahoric and fixes
the pinned frame. For the second take the hyperplane average with
positive exponent in Q87F4. A block lift above the split prime is
paired on the other half with its inverse Hermitian adjoint.
In the matrix chart the test at the inverse lift is left
multiplication by $\operatorname{diag}(\widetilde b,1)$ times the scalar
$\chi(\det b)^{-1}$: the forward lift acts on the half-invariant
functional by $\chi(\det b)$, as follows by evaluating the split
Weil determinant character on its moving block. Its flag level
character is one, since its complementary block is the identity
and every other designated residue is the identity. The factor
$\chi(\det b)$ in the positive average therefore cancels. The principal-level invariance
makes the equality independent of all chosen lifts.

We prove these conditions force $H$ into the span of the right
translates of $\mathbf1_B$. For $X$ put $S=\operatorname{im}_K X$,
$U=S\cap O^{r+t}$ and let $\bar U$ be its reduction. The top
coordinates of the vectors in $\bar U$ with bottom zero form a
subspace $D\subseteq k^r$. If $D=k^r$, lift the $r$ standard
vectors to $U$. Their top blocks are independent over $K$, so $X$
has rank $r$ and they equal the columns of $Xg$ for some
$g\in\mathrm{GL}_r(K)$. In particular $Xg\in B$.

If $D$ is proper, choose $a\ne0$ killing it. For $b\in D_a$ the
linear map $(z,w)\mapsto(1-b)z$ on $\bar U$ factors through its
bottom coordinate $w$, since its restriction to the kernel of that
coordinate is zero. Extend the factor to $k^t$ and call its matrix
$b'$. Choose an integral block upper triangular matrix $Q_0$ with
reduction $\left(\begin{smallmatrix}b&b'\\0&1\end{smallmatrix}\right)$.
Its reduction fixes $\bar U$. Saturation of $U$ makes $O^{r+t}/U$
free over the DVR. Choose an integral projection $P$ onto $U$.
Then $Q=Q_0+(1-Q_0)P$ fixes $S$ pointwise and has the same
invertible reduction as $Q_0$. Since
$\operatorname{diag}(\widetilde b,1)Q^{-1}$ fixes $e$ on reduction,
the first condition makes every term in the second equal to $H(X)$.
The subgroup has positive cardinality, so $H(X)=0$.

On each right translate $B g^{-1}$ the first condition makes $H$
constant. Indeed two matrices reducing to $e$ extend to invertible
integral matrices reducing to the identity by adjoining the last
$t$ standard columns. The ratio gives the required source mover.
Two right translates of $B$ are disjoint or equal: a common point
has invertible top block after either change of frame, hence the
ratio of the two right multipliers belongs to
$1+\varpi M_r(O)$ and fixes $B$. These open sets cover the
nonzero support; local constancy and compact support give a finite
subcover. This proves the claimed finite expansion.

At a twist prime the prescribed function in the conjugate half is
supported on integral $X$ with bottom block zero modulo $\varpi$
and invertible top block $A$ modulo $\varpi$. Its value on this
support is $\xi(\det A)^{-1}$ with the member's nontrivial
residue character $\xi$. Here $H$ satisfies the same inverse
determinant law under the left parahoric preserving the top block
on reduction: apply its level law to $X$ and to $e$, where the
prescribed function has value one. The two Weil scalars agree and
cancel. The first condition above follows. The second follows from
$\sum_{b\in D_a}\xi(\det b)^{-1}=0$, since $\det:D_a\to k^\times$
is onto. Thus $H$ is again spanned by frame indicators. Project this
expansion onto the required determinant character of the
$\mathrm{GL}_r(k)$ block, using normalized averaging. The action
is well defined on principal-type functions and commutes with right
translations. The projection fixes $H$. On $\mathbf1_B$ its support
is the stated twist support and its value at $e$ is
$1/|\mathrm{GL}_r(k)|$, since only the identity residue contributes.
By the determinant law its value everywhere is this same constant
times the prescribed twist test. This gives generation by that
particular test at each twist prime.
% END Q87F6

% BEGIN Q87F7
\paragraph{Q87F7. A nonzero varying-test lift detects the prescribed lift.}
On the original member with the real and level realizations above,
suppose $y$ is good in the sense of $\mathrm{GoodVt}$. If
$I_y(f,\ell,g)\ne0$ for some Schwartz function $f$, coefficient
functional and target argument, then $\mathrm{thG}(y)\ne0$.

Indeed a zero prescribed lift gives $\mathrm{thF}(y)=0$ by Q87F1.
Q87F2--3 give zero for every test with the fixed finite annihilator.
N86V1 and Q87F4 replace any original test by a clean type test
with the same integrals. Use the finite-level decomposition described
in Q87F3. There are only finitely many nonstandard places:
bounded support and local constancy in the finite adeles give this
outside the finitely many designated and twist places. In a basis
of the other-variable coefficients the character and averaging
conditions at a chosen place hold componentwise. Q87F5--6 write
each local coefficient as a finite sum of target translates of
the prescribed local factor. Equivalently, one can expand first
in a basis of the local span; the coefficients in the remaining
variables are linear combinations of the original coefficients
and satisfy all conditions at their own places. The finite-level
groups and residue characters are products, and lifts supported at
one base place exist by the split frames or the self-dual unramified
model. Thus these decompositions can be performed one place at a
time. After the last place, each term is a target translate of a
test with the same finite factor as $v_P$. Such a test satisfies
every finite annihilator relation of $v_P$. The target translation
changes $g$ to $gg'$ and introduces no new test in the vanishing
statement. All the integrals are therefore zero, a contradiction.

In particular this proves the fixed-test part of Y86N2, including
$\mathrm{GenFin}$ with the native unramified lattice theorem.
Applied to every nonzero eigenvector with the required string and
good projector, the assertion $\mathrm{SomeLift}$ implies the
original $\mathrm{TypeNZ}$. This proof of descent uses the same
finite test, coefficient representation and Haar measures as the
original downward map.
% END Q87F7
